\chapter{Ամփոփ ցուցադրման սցենար}

Դիպլոմի վերջնական ցուցադրման հնարավոր սցենար.

\begin{enumerate}[leftmargin=1.6em]
  \item Օգտվողը բացում է Web հարթակը և ընտրում \textbf{Vibration} սենսորը և
        \textbf{Periodic} ազդանշանը։
  \item Կոնֆիգուրացնում է պարամետրերը՝ sample rate, ամպլիտուդ, աղմուկ, threshold,
        filter, anomaly, և սեղմում \textbf{Start Experiment}։
  \item Backend-ը կոնֆիգուրացիան ուղարկում է AUP ZU3-ին. Zynq PS-ը պարամետրերը
        գրում է FPGA-ի AXI4-Lite ռեգիստրներում։
  \item FPGA-ն սկսում է գեներացնել վիբրացիանման տվյալներ և կատարում իրական ժամանակի
        նախնական մշակում։
  \item Տվյալները փոխանցվում են PS, ապա Ethernet-ով՝ սերվեր, և պահվում տվյալների բազայում։
  \item Web dashboard-ը ցուցադրում է raw և processed ազդանշանները իրական ժամանակում։
  \item FPGA-ն ներարկում է անոմալիա. FPGA-ի threshold detector-ը հայտնաբերում է
        իրադարձությունը, իսկ սերվերային մոդելը՝ նույնպես վերլուծում այն։
  \item Հարթակը համեմատում է FPGA-ի հայտնաբերումը և ծրագրային վերլուծությունը։
  \item Օգտվողը Web հարթակից փոխում է FPGA-ի պարամետրերը, և FPGA-ի վարքը փոխվում է
        \textbf{runtime}-ում՝ առանց bitstream վերագեներացնելու։
\end{enumerate}

Վերջնական արդյունքը ուստի ոչ թե պարզ սենսորային dashboard է, այլ ամբողջական,
կոնֆիգուրելի \textbf{Web Platform $\leftrightarrow$ Server $\leftrightarrow$ Zynq PS
$\leftrightarrow$ FPGA} համակարգ, որն աջակցում է տվյալների գեներացում, հավաքագրում,
իրական ժամանակի մշակում, կոնֆիգուրացիա, պահպանում, վիզուալիզացիա, վերլուծություն և
ընտրովի ադապտիվ կառավարում։
